\chapter{Transacciones y crash refinement}
\label{cap:transacciones}

Un filesystem que vive en disco tiene que sobrevivir \emph{crashes}: el
power-cut, el kernel panic, el hot-yank del SSD. La pregunta no es ``¿se
rompe?'' sino ``¿en qué estado queda?''. Hay tres respuestas posibles,
en orden creciente de garantías:

\begin{description}
  \item[\emph{Inconsistent}.] El FS queda corrupto; \texttt{fsck}
    intenta repararlo con heurísticas. Bueno hasta el primer archivo
    crítico que no se recuperó.
  \item[\emph{Eventually consistent}.] Tras el reboot, una secuencia de
    pasos automáticos (replay del journal, reconciliación) llega a un
    estado consistente \emph{eventualmente}. \texttt{ext4} con
    \texttt{data=journal} entra acá.
  \item[\emph{Always consistent}.] El FS está \emph{en todo momento} en
    un estado consistente. Tras el crash, la primera lectura ya retorna
    datos correctos. Es a lo que sotFS apunta.
\end{description}

Este capítulo presenta la herramienta para llegar a la tercera: las
\emph{transacciones de grafo} (GTXN), su modelo formal con
\emph{crash refinement} en TLA\textsuperscript{+}, y su implementación
en el runtime, que difiere honestamente del modelo en un punto
importante.

\section{Modelo formal vs.\ implementación}

\begin{table}[h]
\centering
\renewcommand{\arraystretch}{1.25}
\begin{tabularx}{\linewidth}{l X X}
  \toprule
   & \textbf{Modelo formal (TLA\textsuperscript{+})} & \textbf{Implementación (Rust)} \\
  \midrule
  Spec / módulo & \citaTLA{sotfs\_transactions.tla} + \citaTLA{sotfs\_crash.tla} & \citaCrate{sotfs-tx} sobre \citaCrate{sotfs-storage} \\
  Mecanismo & \wal{} explícito + commit marker & snapshot in-memory + redb txn \\
  Atomicidad & 2-phase commit modelado paso a paso & \texttt{with\_transaction} closure \\
  Durabilidad & \fsync{} sobre \wal & \texttt{redb::Database::commit} (interno: \fsync) \\
  Crash safety & teorema \texttt{Spec\_Concrete \(\Rightarrow\) Spec\_Abstract} & confiada al backend redb \\
  Refinement & \citaTLA{sotfs\_crash\_refinement.tla} & por inspección \\
  \bottomrule
\end{tabularx}
\caption{Modelo formal vs.\ implementación. El modelo describe un \wal{}
explícito; el runtime delega la durabilidad al backend redb. La
correspondencia es por \emph{inspección humana}, no por \emph{extracción
mecánica}; cap.~\ref{cap:tla} discute las consecuencias de esa elección.}
\label{tab:modelo-vs-impl}
\end{table}

Es importante que esto se enuncie en el primer párrafo del capítulo, no
en una nota al pie: el \wal{} de los specs es una construcción
\emph{teórica} para razonar sobre crashes; el runtime no escribe un
\wal{} propio sino que confía en las garantías del backend redb. La
prueba formal \texttt{Spec\_Concrete \(\Rightarrow\) Spec\_Abstract} habla
del \wal{} teórico, no del código de producción. La virtud de la
separación: si mañana cambia el backend (de redb a otro KV embebido), la
prueba formal sigue valiendo siempre que el backend nuevo dé las mismas
garantías de atomicidad y durabilidad.

\section{La GTXN del runtime}

La unidad de trabajo en \citaCrate{sotfs-tx} es la \emph{Graph
Transaction} (GTXN). Su API es minimalista
(\citaLine{sotfs-tx/src/lib.rs}{22}):

\begin{lstlisting}[language=rust, caption={GTXN como wrapper sobre el TypeGraph.}, label={lst:gtxn}]
pub struct Gtxn {
    pub state: GtxnState,
    snapshot: TypeGraph,
}

impl Gtxn {
    pub fn begin(graph: &TypeGraph) -> Self { ... }
    pub fn commit(&mut self) -> Result<(), GraphError> { ... }
    pub fn rollback(self, graph: &mut TypeGraph) { ... }
}

pub fn with_transaction<F, T>(
    graph: &mut TypeGraph, f: F
) -> Result<T, GraphError>
where F: FnOnce(&mut TypeGraph) -> Result<T, GraphError>;
\end{lstlisting}

El patrón canónico es la closure \texttt{with\_transaction}, que
\emph{snapshotea} el grafo al inicio, ejecuta la closure, y o bien
verifica los invariantes y commitea, o rollbackea sobre cualquier error
o violación de invariante.

\begin{algorithm}[H]
\SetAlgoLined
\KwIn{Grafo $G$, closure $f: G \to \textrm{Result}$}
\KwOut{Resultado de $f$ con $G$ commiteado, o $G$ restaurado al snapshot}
\caption{\textsc{WithTransaction}}\label{alg:with-tx}


$\textrm{snap} \gets \textrm{clone}(G)$\;
$r \gets f(G)$\;
\If{$r = \texttt{Ok}(\cdot)$}{
  \If{$G.\textsc{CheckInvariants}() = \texttt{Err}$}{
    $G \gets \textrm{snap}$\;
    \Return{\texttt{Err(invariante violada)}}\;
  }
  
  \Return{$r$}\;
}

$G \gets \textrm{snap}$\;
\Return{$r$}\;
\end{algorithm}

Notas críticas sobre la implementación actual:

\begin{itemize}
  \item El \emph{snapshot} es una copia completa del \texttt{TypeGraph}
    (\texttt{graph.clone()}). Para grafos pequeños es barato; para grafos
    de millones de nodos, el costo es prohibitivo. La estructura RCU
    del cap.~\ref{cap:implementacion} mitiga esto en lectura, pero la
    snapshot de transacción todavía clona.
  \item No hay todavía un mecanismo de \emph{commit en dos fases} entre
    objetos: GTXN es siempre single-host. La extensión Tier-2 (multi-host
    2PC) está en el roadmap (\citaCrate{sotfs-experimental}).
  \item La durabilidad llega cuando el caller invoca
    \texttt{RedbBackend::save(graph)}, no en el commit de la GTXN. La
    GTXN garantiza atomicidad in-memory, no durabilidad on-disk.
\end{itemize}

\section{El protocolo abstracto del modelo TLA: GTXN con WAL}

Aunque el runtime no escriba \wal{} propio, el modelo formal del
\citaTLA{sotfs\_crash.tla} sí lo hace, y razonar sobre crash safety
requiere conocerlo. El protocolo de cuatro fases:

\begin{algorithm}[H]
\SetAlgoLined
\KwIn{Transacción $\textit{tx}$ con conjunto de cambios
$\mathit{tx.changes}$}
\caption{\textsc{GtxnCommit} (modelo formal)}\label{alg:gtxn-commit}


$\textit{rec} \gets \textsc{Serialize}(\mathit{tx.changes})$\;
$\textsc{Disk.Write}(\textit{wal\_log}, \textit{rec})$\;

$\textsc{Disk.Fsync}(\textit{wal\_log})$\;


$\textsc{Disk.Write}(\textit{wal\_log}, \texttt{COMMIT\_MARKER})$\;

$\textsc{Disk.Fsync}(\textit{wal\_log})$\;


\ForEach{cambio $c \in \mathit{tx.changes}$}{
  $\textsc{Disk.Write}(\textit{graph\_store}, c)$\;
}
$\textsc{Disk.Fsync}(\textit{graph\_store})$\;


$\textsc{Wal.Truncate}(\textit{tx})$\;
\end{algorithm}

El \emph{commit point} es el segundo \fsync{} (línea 5).
\emph{Antes} del commit point, el crash deja \wal{} parcial sin marker
$\Rightarrow$ el recovery descarta las operaciones (vuelve al estado
pre-tx). \emph{Después} del commit point, el crash deja \wal{} con
marker $\Rightarrow$ el recovery re-aplica las operaciones (continúa al
estado post-tx). Nunca un estado intermedio.

\section{Crash refinement}

El spec abstracto \citaTLA{sotfs\_transactions.tla} modela GTXN
\emph{sin} crashes; el spec concreto \citaTLA{sotfs\_crash.tla}
\emph{con} crashes. La técnica de \emph{crash refinement}
\cite{Sigurbjarnarson2016} prueba que toda traza concreta es traza
abstracta vía un \emph{refinement mapping}.

\begin{lstlisting}[language=tla, caption={Acciones \texttt{Crash} y
\texttt{Recover} en \citaTLA{sotfs\_crash.tla} (extracto).}, label={lst:tla-crash}]
\* Crash puede ocurrir en cualquier momento.
Crash ==
    /\ ~crashed
    /\ crashed' = TRUE
    /\ memState' = NULL                  \* memoria volátil borrada
    /\ UNCHANGED <<diskState, walState, gtxnState>>

\* Recovery: replay del WAL desde el último commit marker.
Recover ==
    /\ crashed
    /\ memState' = ReplayWAL(walState, diskState)
    /\ crashed' = FALSE
    /\ UNCHANGED <<diskState, walState>>
\end{lstlisting}

\subsection*{El refinement mapping}

\begin{lstlisting}[language=tla, caption={Mapping de refinement
concreto $\to$ abstracto.}, label={lst:tla-refinement}]
RefinementMapping ==
    LET
        ConcreteState == <<diskState, walState, memState>>
        AbstractState ==
            IF crashed THEN
                \* mapeamos al estado abstracto "como si la última
                \* committed tx fuera la actual"
                ReplayWAL(walState, diskState)
            ELSE
                memState
    IN
        \A op \in TxOps :
            ConcreteOp(op) => AbstractOp(op)

THEOREM Spec_Concrete => Spec_Abstract
\end{lstlisting}

Bajo TLC con configuración \texttt{small}, el teorema se verifica
exhaustivamente sobre estados acotados (cap.~\ref{cap:tla} lo desarrolla).

\section{Las propiedades verificadas}

\begin{description}
  \item[\textbf{No partial application.}] Tras \texttt{Recover},
    \emph{toda} operación de una GTXN está aplicada o \emph{ninguna}.
    Nunca ``se aplicaron 3 de 5''.
  \item[\textbf{Durability.}] Si \texttt{gtxn\_commit} retornó
    \texttt{Ok}, la transacción sobrevive cualquier crash posterior.
  \item[\textbf{Recovery consistency.}] Tras \texttt{Recover}, el estado
    de memoria satisface \emph{todas} las invariantes del
    cap.~\ref{cap:typegraph} (los siete + el de índice).
  \item[\textbf{WAL integrity.}] Una vez escrito el commit marker, el
    contenido del \wal{} no se modifica hasta el GC; el GC sólo borra
    entradas \emph{ya aplicadas}.
\end{description}

\section{El protocolo de fsync}

La barrera de durabilidad es \fsync: pide al SSD que flushee sus cachés
volátiles a flash NAND. Sin \fsync, una escritura puede quedarse en el
caché del controlador y perderse en un power-cut.

\begin{lstlisting}[language=shellcmd, caption={Pipeline conceptual de
\fsync.}]
disk.write(buf)         # bytes a kernel page cache
disk.fsync()            # kernel page cache -> driver buffer
                        # driver buffer -> SSD controller
                        # SSD controller -> NAND
                        # ack al kernel
\end{lstlisting}

\begin{invariante}[Persistencia tras \fsync]
\label{inv:fsync-persist}
Tras \fsync\texttt{(loc)} retornar exitosamente, los bytes escritos antes
de la llamada a \texttt{loc} están persistidos en NAND y sobrevivirán
cualquier power-cut posterior.
\end{invariante}

La invariante depende del SSD respetar el comando \texttt{FUA} (\emph{Force
Unit Access}). SSDs consumer-grade malos lo ignoran; SSDs enterprise-grade
lo respetan. Verificar esto en runtime con un test de power-cut físico es
la única vía; el modelo asume \fsync{} honesto.

\section{Trampas históricas}

\begin{trampa}[\texttt{write A; write B; fsync} no garantiza orden]
\label{tr:write-reorder}
Versión vulnerable: el código asume que dos \texttt{write} seguidos por
un \fsync{} se persisten en orden $A, B$. \emph{No es así}. El SSD puede
reordenar internamente; \fsync{} sólo garantiza que \emph{ambos} están
persistidos al retornar, no el orden.

Si el código depende de que $A$ esté en disco antes que $B$ (e.g.\ $B$
es el commit marker que valida $A$), un crash entre los dos \fsync{}
puede dejar $B$ en disco sin $A$.

\textbf{Fix.} \texttt{write A; fsync; write B; fsync}. Dos \fsync{}
explícitos. Costo: latencia, especialmente en HDD. Justificable cuando
el orden importa (commit marker), no cuando no (datos paralelos).

\textbf{Lección.} \fsync{} es una barrera de \emph{durabilidad puntual},
no de \emph{ordering entre operaciones}. Cada barrier tiene su propio
\fsync.
\end{trampa}

\begin{trampa}[GC del WAL antes del fsync del graph]
\label{tr:wal-truncate-early}
\citaCommit{4775904}. Versión vulnerable: el orden de Fase 3 y Fase 4
estaba invertido. El \wal{} se truncaba (Fase 4) \emph{antes} de que el
graph store hiciera \fsync{} (Fase 3 incompleta).

Si crash ocurre después del \wal{} truncate pero antes del graph
\fsync: el \wal{} no tiene la transacción (truncado), el graph no la
tiene (no \fsync-eado). \emph{Pérdida silenciosa}.

\textbf{Fix.} El orden estricto Fase 1 $\to$ Fase 2 $\to$ Fase 3 $\to$
Fase 4. Cada fase termina con su \fsync{} correspondiente. \wal{}
truncate sólo tras graph \fsync{} exitoso.

\textbf{Lección.} La durabilidad es un pipeline ordenado de barreras.
Una operación de \emph{cleanup} (truncate del \wal) DEBE preceder de
todas las operaciones que dependen de los datos cleaned-up. Caer en la
tentación de hacer cleanup \emph{en paralelo} para ``ahorrar latencia''
rompe la cadena de durabilidad.
\end{trampa}

\section{La invariante de \emph{readers see committed-or-pre-tx}}

\begin{invariante}[Read isolation]
\label{inv:read-iso}
Un \texttt{read} concurrente con un \texttt{gtxn\_commit} en progreso ve
\emph{o} el estado anterior a la transacción \emph{o} el estado posterior;
nunca un estado intermedio. La elección depende de cuándo se linealiza el
read respecto del commit point.
\end{invariante}

Implementación: los readers usan el Tier 0 \rcu{}
(cap.~\ref{cap:implementacion}); toman snapshot de la \emph{época
global} del FS al iniciar la lectura, leen consistentemente de esa
snapshot. Los writers (gtxn) bumpean la época sólo después del commit
point, de modo que los readers con snapshot vieja siguen viendo el
estado pre-tx, los nuevos ven el post-tx.

\section{Origen del diseño}

\wal{} viene de bases de datos (\textsc{Aries} \cite{Mohan1992}); su
aplicación a filesystems es de \texttt{ext3} y \texttt{XFS}. La idea de
\emph{crash refinement} (modelo concreto con crashes refinando un modelo
abstracto sin) es de Yggdrasil \cite{Sigurbjarnarson2016}. La
combinación con un grafo tipado es contribución de sotFS; la dificultad
es mantener la \emph{inmutabilidad estructural} del grafo bajo crashes
(los nodos viejos sobreviven, los nuevos pueden estar parciales).

\section{Resumen del capítulo}

\begin{itemize}
  \item GTXN protocoliza la mutación durable en el modelo formal:
    \wal{} $\to$ commit marker $\to$ graph apply $\to$ \wal{} GC, con
    \fsync{} en cada paso.
  \item En el runtime v0.2.5 la GTXN es snapshot-based in-memory;
    la durabilidad se delega al backend redb.
  \item El \emph{commit point} es el \fsync{} del marker en el modelo;
    antes, crash descarta; después, crash re-aplica.
  \item Crash refinement (\citaTLA{sotfs\_crash\_refinement.tla})
    prueba que toda traza concreta con crashes es traza abstracta sin
    ellos.
  \item Read isolation: los readers usan Tier 0 \rcu{} sobre la época
    FS; ven o pre-tx o post-tx, nunca intermedio.
  \item Trampas: write reordering sin \fsync{} intermedio rompe la
    durabilidad ordenada; \wal{} truncate antes del graph \fsync{}
    causa pérdida silenciosa.
\end{itemize}

\section*{Ejercicios}

\begin{ejercicio}[\dificultad{$\star$}]
\label{ej:tx-1}
Diga el orden \emph{correcto} de los cuatro \fsync{} en una GTXN del
modelo formal. ¿Cuál es el commit point? ¿Qué pasa si un crash ocurre
entre el \fsync{} 1 y el \fsync{} 2?
\end{ejercicio}

\begin{ejercicio}[\dificultad{$\star$}]
\label{ej:tx-2}
Enuncie en una línea cada una de las cuatro propiedades verificadas en
sec.~``Las propiedades verificadas''. ¿Cuál es la \emph{más difícil} de
demostrar? Justifique brevemente.
\end{ejercicio}

\begin{ejercicio}[\dificultad{$\star\star$}]
\label{ej:tx-3}
Lea la propiedad \emph{Recovery consistency}. ¿Qué propiedades del
cap.~\ref{cap:typegraph} están en juego? ¿Bastaría verificar sólo el
\invref{inv:link-count} para garantizar consistencia?
\end{ejercicio}

\begin{ejercicio}[\dificultad{$\star\star$}]
\label{ej:tx-4}
La GTXN del runtime usa \texttt{graph.clone()} como snapshot. Estime el
costo en bytes y en tiempo para un grafo de \(10^6\) inodes. ¿En qué
operaciones se nota más?
\end{ejercicio}

\begin{ejercicio}[\dificultad{$\star\star$}]
\label{ej:tx-5}
Un atacante con acceso físico al SSD puede inyectar power-cuts
controlados. ¿Puede forzar pérdida de datos en una GTXN correctamente
implementada (modelo formal)? Justifique en términos del
\invref{inv:fsync-persist} y los assumptions sobre el SSD.
\end{ejercicio}

\begin{ejercicio}[\dificultad{$\star\star$}]
\label{ej:tx-6}
La \trampref{tr:wal-truncate-early} se cierra con un orden estricto.
Explique por qué un \emph{commit en dos fases} entre el \wal{} y el
graph store es necesario aún cuando ambos viven en el mismo SSD.
\end{ejercicio}

\begin{ejercicio}[\dificultad{$\star\star$}]
\label{ej:tx-7}
Lea \citaLine{sotfs-tx/src/lib.rs}{55} (\texttt{with\_transaction}).
¿Qué pasa si la closure $f$ commitea el grafo en estado consistente pero
\emph{no} llamó a \texttt{check\_invariants}? Pista: revisar el orden de
las líneas.
\end{ejercicio}

\begin{ejercicio}[\dificultad{$\star\star$}]
\label{ej:tx-8}
Tabla~\ref{tab:modelo-vs-impl}: explique por qué la fila ``Durabilidad''
no es estrictamente una correspondencia (el runtime no escribe \wal{}
propio). ¿Qué garantías del backend redb son necesarias para que el
runtime herede las propiedades del modelo?
\end{ejercicio}

\begin{ejercicio}[\dificultad{$\star\star$}]
\label{ej:tx-9}
Enuncie la \invref{inv:read-iso} sin usar la palabra ``snapshot''.
Pista: hablar de visibilidad y orden lineal de eventos.
\end{ejercicio}

\begin{ejercicio}[\dificultad{$\star\star\star$}]
\label{ej:tx-10}
Diseñe un test de \emph{crash injection}: una herramienta que mata el
proceso del runtime FUSE en momentos aleatorios y verifica que el FS está
en un estado consistente tras boot. Pista: \texttt{fork(2)} + \texttt{kill -9}
hijo a tiempo random + script de assertion.
\end{ejercicio}

\begin{ejercicio}[\dificultad{$\star\star\star$}]
\label{ej:tx-11}
Compare las garantías de crash safety de sotFS contra \texttt{ext4}
(con \texttt{data=journal}), \texttt{btrfs} y ZFS. ¿Qué ofrece sotFS
que los otros no? ¿Qué ofrecen los otros que sotFS no?
\end{ejercicio}

\begin{ejercicio}[\dificultad{$\star\star\star$}]
\label{ej:tx-12}
Diseñe la extensión Tier-2 (multi-host 2PC) de GTXN. ¿Qué nodos del
TypeGraph hacen falta agregar (un nodo coordinator)? ¿Qué aristas? ¿Qué
spec TLA\textsuperscript{+} hay que extender?
\end{ejercicio}
